\SourceLeaf{160}
\textit{reine} Gewinn der Familie durch die Aufnahme in die Kasse (nicht wie er wirklich gewesen ist, sondern wie er bei jetziger Höhe der Beiträge und Pensionen gewesen sein würde) zu 4999 Rthl. 4 Ggr. Werden aber jene Summen durch Discontirung à 4 Proc. auf den Zeitpunkt des Beitritts, Michaelis 1763, zurückgeführt, so reduciren sich
\begin{center}
\begin{tabular}{l l r l r l}
die Beiträge&auf&211&Rthl.&23&Ggr.\\
die Pensionsbezüge&auf&527&,,&10&,,
\end{tabular}
\end{center}
\noindent und mithin jener grosse Gewinn, nach seinem Geldwerth für diese Epoche, auf die geringe Summe von 315 Rthl. 11 Ggr. Eine Rechnungsprobe würde haarscharf ergeben, dass dieses kleine Kapital, jährlich durch die Zulegung von 10 Rthl. und durch die Zinsen zu 4 Proc. vermehrt, und so fortwährend sich vergrössernd, am Schluss des Jahres 1812 eine solche Grösse erlangt haben würde, dass von da an durch die fernern Zinsen und theilweise durch Einziehen des Kapitals die Pension genau bis Ende November 1834 gedeckt war.

Aus diesem Grunde, so wie noch aus andern, deren Entwicklung mich hier zu weit führen würde, ist der mittlere Geldwerth der Aufnahme in die Gesellschaft, in dem Sinn wie er hier verstanden wird, gar nicht vergleichbar mit dem durchschnittlichen Geldwerth der Ansprüche aller in einem bestimmten Zeitpunkte coexistirenden Mitglieder, von denen immer ein grosser Theil schon sehr lange in der Gesellschaft gewesen ist.

Träte nun jedes Jahr \textit{Ein} neues Mitglied ein, dessen Berechtigung durchschnittlich zu obigem Geldwerth angeschlagen werden müsste, so ist klar, dass zur Deckung dieser sich immer neu gebärenden Ansprüche die Zinsen eines eisernen Kapitals erforderlich und zureichend sein würden, dessen Höhe
\begin{center}
\begin{tabular}{l}
bei dem Zinsfusse von \(3\tfrac{1}{2}\) Proc. zu 17182,86Thalern\\
bei dem Zinsfusse von 4 Proc. zu 12912,50 Thalern
\end{tabular}
\end{center}
\noindent angenommen werden müsste.

Bis hieher ruhet die Rechnung auf einer sichern Grundlage. Es fehlt nun aber zur Vollendung derselben noch ein wesentliches Element, für welches jede im Voraus gewagte Schätzung nach der Natur der Sache nur eine sehr precäre Geltung haben kann, ich meine die Durchschnittszahl der künftig jährlich neu beitretenden Theilnehmer. Befragt man die Erfahrung, so zeigt sich in der neuern Zeit eine sehr starke Vergrösserung dieses Elements. Für die ersten achtzig Jahre des Bestehens der Anstalt kann man die Durchschnittszahl der jährlich beitretenden nahe zu \(1\tfrac{1}{2}\) annehmen; allein seit den letzten 20 Jahren sind überhaupt 64 Beitritte\footnote{Unter denselben ist G. mitgezählt, dessen Witwe Pension geniesst, der aber selbst eigentlich nicht recipirt gewesen war. Wenn man aber voraussetzen muss, dass in einem ähnlichen künftigen Falle auch wieder abseiten der Verwaltung ein ähnlicher Beschluss gefasst werden wird, so dürfen solche, wenn auch seltene, Fälle nicht unberücksichtigt bleiben, sondern müssen in einer längern Reihe von Erfahrungen als \textit{auch} mögliche Chancen mitgezählt werden.} erfolgt, wovon 23 auf das Decennium 1825—1835, und 41 auf das Decennium 1835—1845 kommen. Das Mittel aus den letzten 20 Jahren wäre demnach \(3\tfrac{1}{5}\), oder, wenn man G. nicht mitzählen will, \(3\tfrac{3}{20}\). Für die letzten 10 Jahre allein ist dies Mittel \(4\tfrac{1}{10}\), und für die letzten 5 Jahre allein, während welcher 27 neue Mitglieder eingetreten sind, sogar \(5\tfrac{2}{5}\). Man erkennt hieraus, ein wie missliches Unternehmen es ist, hienach ein Prognosticon für die Zukunft zu stellen. Aber geringer als \(2\tfrac{1}{2}\) würde ich doch nicht wagen, die künftige jährliche Durchschnittszahl anzunehmen, und dieses Verhältniss würde auch wie ganz harmonisch zu der jetzigen Interessentenzahl 51 betrachtet werden können, insofern die bisherigen Erfahrungen 20 Jahre oder etwas weniges mehr als durchschnittliche Dauer der Theilnahme herausgestellt haben. In dieser Hypothese stellt sich also der gegenwärtige Geldwerth
